%=============================================================================!
% 17.5  LiC6 WITH A FOUNDATION MODEL                                          !
%=============================================================================!
\section{\texorpdfstring{LiC$_6$}{LiC6} with a foundation model}
\label{sec:ip-lic6}

Books pushed between the shelves of a bookcase whose shelves can slide
force them apart, while the shelves themselves keep their length. Lithium
does this to graphite. Chapter~14 could only imitate it with a model solid
(\cref{sec:pr-shape}); with MACE-MP-0 and D3 the real material can be
run, and this section does so for the fully lithiated compound LiC$_6$ at
room temperature, the study that Chapters~13 and~14 left for a real
potential. Its numbers are the model's predictions (\cref{sec:ip-aimd}).

\subsection*{Building the cells}

ASE builds the structures. Graphite's cell holds four atoms, two in each
of its sheets stacked A, B, with Trucano and Chen's lattice constants as
the start. In LiC$_6$ the sheets stack A, A and a lithium atom sits over
every third hexagon, the pattern written
$\sqrt3\times\sqrt3$\,R30$^\circ$: its cell is the sheet's own cell made
$\sqrt3$ times longer and turned by \ang{30}, built in ASE as the
supercell of the two-atom cell of graphene, a single sheet of graphite,
whose vectors $\vb{a}$ and $\vb{b}$ have length $a$ and lie \ang{120}
apart, with the vectors $2\vb{a} + \vb{b}$ and $-\vb{a} + \vb{b}$
(\cref{ex:ip-supercell}), and it holds six carbon atoms and one lithium.
Relaxed together with its cell at \SI{0}{\kelvin}, to forces below
\SI{1}{\milli\electronvolt\per\angstrom} and stresses below
\SI{0.02}{\giga\pascal}, LiC$_6$'s sheets lie \SI{3.616}{\angstrom} apart,
2.3\% below the measured \SI{3.70}{\angstrom} quoted by Hazrati, de Wijs
and Brocks\cite{Hazrati2014}, and the sheet's lattice constant $a$, the
side of its two-atom cell, is \SI{2.4864}{\angstrom}, 0.88\% longer than
the \SI{2.4647}{\angstrom} of graphite relaxed with the same model: the
sheets stretch a little, as the electrons that lithium gives them are
expected to make them do.

\subsection*{Anisotropic and isotropic coupling at \SI{300}{\kelvin}}

A cell of $3\times3\times1$ of these units, 63 atoms, was run at
\SI{300}{\kelvin} and zero pressure with ASE's \texttt{MaskedMTKNPT}, which
lets each lattice vector change its length while keeping its direction,
with $\tau_{\mathrm T} = \SI{100}{\femto\second}$ and $\tau_{\mathrm P} =
\SI{1}{\pico\second}$, for \SI{20}{\pico\second}; and with
\texttt{IsotropicMTKNPT}, which scales the three lengths together, for
\SI{10}{\pico\second} (\cref{fig:ip-lic6}). The isotropic run keeps the
ratio of the lengths it started from, and the stiff sheets hold the
in-plane length, and with it the spacing, at their values at
\SI{0}{\kelvin}: the spacing stays at $(3.6162 \pm
0.0005)$\,\si{\angstrom}. The diagonal of the pressure tensor, computed
from the model's stress with the atoms' kinetic part added, on 25 frames
of the last \SI{5}{\pico\second}, is $(-1.00 \pm 0.45)$, $(-1.11 \pm
0.36)$ and $(1.68 \pm 0.10)$\,\si{\giga\pascal} along $x$, $y$ and $z$:
the sheets push apart along $z$ while each sheet pulls in along $x$ and
$y$, and only the mean is zero. With the lengths coupled apart, all
three components, from 75 frames of the last \SI{15}{\pico\second}, are
zero within 1.2 of their errors, $(0.06 \pm 0.28)$, $(-0.21 \pm 0.18)$
and $(0.01 \pm 0.06)$\,\si{\giga\pascal}, and the spacing opens on
warming.

\subsection*{How long the spacing takes}

The in-plane length settles at once: over the \SI{20}{\pico\second} the
lattice constant $a$ is $(2.4836 \pm 0.0004)$\,\si{\angstrom}, its
quarters agreeing to \SI{0.001}{\angstrom}; \texttt{convergence\_report}
flags it only because its blocks give a smaller error,
\SI{0.0002}{\angstrom}, than the one quoted. The spacing does not settle. Its
means over the four quarters of the run are 3.756, 3.760, 3.766 and
\SI{3.778}{\angstrom}, still rising, and \texttt{convergence\_report}
refuses it with 82 independent samples, fewer than the 100 that the
criterion of \cref{sec:cv-protocol} asks. A layered solid's soft
direction relaxes slowly. These saved results do not yet establish its
equilibrium spacing; a longer run is required.

\begin{figure}[htb]
\centering
\includegraphics{ch17_practice/lic6}
\caption{LiC$_6$, 63 atoms, at \SI{300}{\kelvin} and zero pressure with
MACE-MP-0 and D3. (a) The spacing of the sheets against time with the
lattice lengths coupled apart (teal) and together (ochre), with the
relaxed spacing at \SI{0}{\kelvin} (grey dotted) and the measured
\SI{3.70}{\angstrom} (red dotted). (b) The running mean of the spacing in
the initial \SI{20}{\pico\second} from the start that
\texttt{convergence\_report} finds,
with twice its error (band). (c) The diagonal of the pressure tensor
under each coupling, with its standard error.}
\label{fig:ip-lic6}
\end{figure}

\begin{keyidea}
A foundation model with D3 relaxes LiC$_6$ to within a few per cent of
its measured spacing, and runs it at room temperature. A layered solid
needs its lattice lengths coupled apart: coupled together, the spacing is
held and the diagonal of the pressure tensor shows the strain. Its soft
direction relaxes over tens of picoseconds, and the convergence report
decides when it has settled.
\end{keyidea}

\notebook{17}{couple LiC6's lengths together and apart, and watch the spacing}

\begin{selfcheck}
In LiC$_6$, how many hexagons of a sheet are there for each lithium atom,
and why does that give six carbon atoms?
\end{selfcheck}
